\scanpage{199}
\mspage{55}
Explicitons le morphisme d'adjonction

\[
 i\longrightarrow u_*u^{\unsure{\wedge|*}}i
\]

(pour $i\in\operatorname{Ob}I$) par ses valeurs sur $j$

\[
 \underset{\substack{\shortparallel\\\operatorname{Hom}(j,i)}}{i(j)}
 \longrightarrow\Gamma(I'_{/j},u^*i) \tag*{(\unsure{x|\ast})}
\]

Soit donc

\[
 u_j\colon I'_{/j}\longrightarrow I,
\]

on a

\[
 \begin{aligned}
 \Gamma(I'_{/j},\underbrace{u^*(i)\deleted{\unsure{u}}\vert I'_{/j}}_{u_j^*(i)})
 &\simeq\operatorname{Hom}(e_{I'_{/j}},u^*(i))\\
 &\overset{\mathrm{adj}}{\simeq}\operatorname{Hom}(u_{j!}e_{I'_{/j}},i),
 \end{aligned}
\]

et la flèche $(x)$, \inserted{fonctorielle en $i$,} est déduite d'une flèche

\[
 u_{j!}(e_{I'_{/j}})\longrightarrow j. \tag*{(\unsure{xx|\ast\ast})}
\]

On voit tout de suite, comme ci-dessus, que $(xx)$ est iso sss pour $\forall\ \alpha\colon k\mathbin{\unsure{-|\to}}j$, \deleted{la flèche} la catégorie $\,{}_{k\backslash}\underset{\alpha}{I'}_{/j}$ est $0$-connexe, \deleted{\unsure{Ma}} et si on l'exige pour $\forall\ j,k,\alpha$, c'est la condition $d''$. Mais il n'est pas clair que pour qu'on ait $(xx)$ iso, il suffise que le $\operatorname{Hom}$ des deux mbs. \unsure{ci-dessus|ci-contre} dans $\forall\ i$ de $I$ (et non \deleted{\unsure{dans}} $F$ dans $I^\wedge$) soit iso. \unsure{--} La difficulté provient du fait que $u_{j!}(e_{I'_{/j}})$ n'est pas forcément représentable, et qu'un \deleted{foncteur} $F$ quelconque de $I^\wedge$
% archive sheet number: 199
